\begin{abstract}
Direct Preference Optimization (DPO) has emerged as an efficient alternative to PPO-based RLHF, yet its application to knowledge-intensive generation tasks reveals a fundamental limitation: standard preference signals---whether from human annotators or LLM judges---exhibit a systematic \textit{verbosity bias} that rewards linguistic fluency over logical correctness. We show empirically that this reward signal blindspot leaves a large logical alignment gap: SFT-trained models achieve NLI entailment scores of only 0.05--0.22, despite producing fluent, confident-sounding text. We propose \textbf{RLearner-LLM}, a framework that resolves this through \textbf{Hybrid-DPO}: an automated preference construction pipeline that fuses a DeBERTa-v3 NLI entailment signal with a verifier LLM score, eliminating the need for human annotation while overcoming the ``alignment tax'' that degrades fluency under single-signal optimization. Evaluated across five academic domains (Biology, Medicine, Law) with three base architectures (LLaMA-2-13B, Qwen3-8B, and Gemma 4 E4B-it), RLearner-LLM achieves up to 6$\times$ NLI improvement over SFT baselines, with NLI gains in $11$ of $15$ (architecture, domain) cells and consistent answer-coverage gains. On Gemma 4 E4B-it---a 4.5B-effective-parameter dense model---Hybrid-DPO lifts NLI entailment in four of five domains (from +11.9\% on Auckland Law to +2.4$\times$ on UK Medicine Year 2), with faster inference across all five, demonstrating that the approach scales down to compact modern base models without losing the alignment-tax mitigation. We further show these gains are robust to aggressive, response-grounded question-quality filtering: under a \textit{strict-discriminator} filter that keeps only questions whose author key matches the modal student answer (computed from per-option student response logs), answer-coverage improves in every architecture--corpus cell and NLI in four of six. Our Qwen3-8B RLearner-LLM wins $95\%$ of blind pairwise comparisons against its own SFT baseline, while GPT-4o-mini in turn wins $95\%$ against our concise RLearner-LLM output---a result that, viewed alongside the $69\%$ win rate the same judge awards a verbose SFT baseline over our DPO-aligned model, reproduces the verbosity bias we document on a frontier comparator and reinforces the case for logic-aware automatic metrics (NLI, ACR) over LLM-as-a-judge on knowledge-intensive generation.\footnote{Code, training scripts, and per-domain evaluation artefacts: \url{https://github.com/14H034160212/Explanation-Generation}}
\end{abstract}
